\section{Introduction}
\IEEEPARstart{T}{he} rapid expansion of complex network infrastructures and ubiquitous IoT ecosystems has precipitated a corresponding escalation in sophisticated cyber-attacks. At the forefront of defending these digital perimeters are Intrusion Detection Systems (IDS), which monitor network traffic to identify malicious activities and policy violations. In recent years, data-driven approaches leveraging Machine Learning (ML) and Deep Learning (DL) have superseded traditional signature-based systems due to their ability to detect zero-day vulnerabilities. However, the transition from theoretical models to real-world deployment reveals a fundamental pathology: network traffic is inherently and extremely imbalanced. Benign packets govern the vast majority of network flows, whereas malicious intrusions—particularly stealthy, targeted attacks like Infiltration or Heartbleed—constitute a microscopic fraction of the data. 

\subsection{Problem Formulation (The Imbalance Dilemma)}
This extreme data asymmetry poses a severe challenge for contemporary anomaly detection. Solitary machine learning models trained on imbalanced datasets naturally drift toward maximizing overall accuracy, indiscriminately classifying minority attacks as benign traffic. This algorithmic bias manifests in two fatal operational flaws: an unacceptably high False Positive (FP) rate, which causes "alert fatigue" among security analysts, and profoundly low recall for critical minority classes, leaving the network vulnerable to devastating breaches. While state-of-the-art algorithms have shown promise, isolated models often lack the robust variance reduction and complex feature mapping required to effectively straddle the decision boundaries of heavily skewed multi-class datasets.

\subsection{Proposed Solution}
To remediate the deficiencies of solitary architectures and raw imbalanced datasets, this paper proposes a comprehensively optimized IDS framework. We engineer a dual-phase solution that tackles both the data-level anomaly and the algorithmic formulation. First, we apply advanced balanced sampling techniques tailored to the CICIDS2017 benchmark dataset, creating a cohesive, balanced distribution that amplifies the signal of minority classes without discarding critical characteristics of the majority traffic. Second, we design and implement a novel hybrid ensemble model. We systematically integrate the discrete, computational efficiency of XGBoost and LightGBM—both renowned for handling high-dimensional, tabular data structures with granular precision—as our primary feature extractors and base learners. To resolve the final classification, we employ a Multi-Layer Perceptron (MLP) as a meta-classifier, effectively mapping the non-linear outputs of the boosting algorithms into a highly accurate, generalized decision matrix.

\subsection{Main Contributions}
The main contributions of this paper are summarized as follows:
\begin{enumerate}
    \item \textbf{Identification and Mitigation of Extreme Imbalance:} We provide a rigorous statistical analysis of the CICIDS2017 dataset and formulate an advanced balanced sampling pipeline to rectify severe class disparities, ensuring minority attack vectors are adequately represented.
    \item \textbf{Novel Hybrid Ensemble Architecture:} We propose a state-of-the-art stacking ensemble that strategically combines XGBoost and LightGBM with an MLP meta-classifier. This specific triangulation exploits both the structural pattern recognition of gradient-boosted trees and the deep associative power of neural networks.
    \item \textbf{Drastic Performance Optimization:} Extensive empirical evaluations demonstrate that our framework dramatically reduces False Positives and elevates the recall rates for historically difficult-to-detect minority attacks, substantially outperforming existing solitary and traditional ensemble approaches.
\end{enumerate}

\subsection{Structure of the Paper}
The remainder of this paper is organized as follows: Section II reviews related literature and current gaps. Section III details the methodology, including data preprocessing, sampling, and the architectural design of the hybrid ensemble. Section IV outlines the experimental setup, while Section V presents an extensive analysis of the results. Finally, Section VI concludes the paper and delineates future research vectors.
